\documentclass[11pt,a4paper]{article}
\usepackage[margin=2.5cm]{geometry}
\usepackage{amsmath,amssymb}
\usepackage{booktabs}
\usepackage[hidelinks]{hyperref}

\title{The three-equation melt condition at an ice--water interface\\[0.3em]
\large and its form in \texttt{solver2d}}
\date{}

\begin{document}
\maketitle

\section{Geometry and notation}

The ice occupies $y<0$ and the water $0<y<L_y$; the ice--water interface is the wall $y=0$ (in
\texttt{solver2d}: \texttt{bc\_bot = 'noslip'}, \texttt{bc\_top = 'freeslip'}, gravity along $+y$).
The unit normal $\mathbf{n}=\mathbf{e}_y$ points from the interface into the water. The subscript $b$
denotes values at the interface ($y=0$). The melt (ablation) rate $m$ is the speed at which the
interface recedes into the ice (in m\,s$^{-1}$ of ice), positive for melting. The effect of pressure
on the freezing point is neglected throughout (it would add a term $\lambda_3 p_b$, with
$\lambda_3\approx -7.5\times10^{-4}$\,$^\circ$C\,dbar$^{-1}$, to \eqref{eq:liquidus}).

\begin{center}
\begin{tabular}{lll}
\toprule
symbol & meaning & typical value \\
\midrule
$T,\ S$ & water temperature, salinity & \\
$T_b,\ S_b$ & temperature and salinity at the interface & unknown \\
$m$ & melt rate (ice thickness per unit time) & unknown \\
$\lambda_1$ & liquidus slope & $\approx -5.7\times10^{-2}$ $^\circ$C\,psu$^{-1}$ \\
$\lambda_2$ & liquidus offset & $\approx 8.3\times10^{-2}$ $^\circ$C \\
$L_f$ & latent heat of fusion & $3.34\times10^{5}$ J\,kg$^{-1}$ \\
$c_w,\ c_i$ & specific heat of water, ice & $\approx 3974$, $2009$ J\,kg$^{-1}$K$^{-1}$ \\
$\rho_w,\ \rho_i$ & density of water, ice & $\approx 1028$, $917$ kg\,m$^{-3}$ \\
$\kappa_T,\ \kappa_S$ & molecular diffusivities of heat and salt in water & $\approx 1.4\times10^{-7}$, $\approx 7\times10^{-10}$ m$^2$s$^{-1}$ \\
$\kappa_i$ & thermal diffusivity of ice & $\approx 1.1\times10^{-6}$ m$^2$s$^{-1}$ \\
$S_i$ & salinity of the ice & $\approx 0$ \\
\bottomrule
\end{tabular}
\end{center}

The liquidus constants differ slightly between sources; the values above are representative of those
used in ice--ocean models (e.g.\ Holland \& Jenkins 1999). Check them, and the other material
properties, against the source you use.

\section{The three equations}

\paragraph{1. Freezing point (liquidus).} The interface is at the local freezing temperature:
\begin{equation}
  T_b = \lambda_1 S_b + \lambda_2 . \label{eq:liquidus}
\end{equation}

\paragraph{2. Heat balance.} The heat conducted from the water to the interface melts ice and is
partly conducted into the ice:
\begin{equation}
  \rho_w c_w \kappa_T \left.\frac{\partial T}{\partial y}\right|_b
  = \rho_i L_f\, m \;+\; \rho_i c_i \kappa_i \left.\frac{\partial T_i}{\partial y}\right|_b . \label{eq:heat}
\end{equation}
The last term (conduction into the ice, $T_i$ the ice temperature) is often neglected, or modelled as
$\rho_i c_i \kappa_i (T_b - T_s)/h_i$ with $T_s$ the temperature at the far side of an ice layer of
thickness $h_i$.

\paragraph{3. Salt balance.} The ice holds (almost) no salt. The salt that the meltwater dilutes at the
interface must be supplied by diffusion in the water:
\begin{equation}
  \rho_w \kappa_S \left.\frac{\partial S}{\partial y}\right|_b = \rho_i\, m\, (S_b - S_i)
  \;\approx\; \rho_i\, m\, S_b . \label{eq:salt}
\end{equation}

Equations \eqref{eq:liquidus}--\eqref{eq:salt} determine the three unknowns $T_b$, $S_b$ and $m$.
For melting ($m>0$) the interface is fresher than the water ($\partial S/\partial y|_b>0$) and warmer
than the freezing point of the far-field water. Since $\kappa_T/\kappa_S \approx 10^2$, salt
diffusion limits the melt rate together with the heat supply.

\paragraph{Signs.} With $\mathbf{n}=\mathbf{e}_y$ pointing into the water, the conductive fluxes
\emph{towards} the interface are $\rho_w c_w \kappa_T\, \partial T/\partial y|_b$ and
$\rho_w \kappa_S\, \partial S/\partial y|_b$. This holds for the set-up above: ice at $y=0$, water at
$y>0$, gravity along $+y$ (the ice is physically on top). If instead the ice is at $y=L_y$ (water at
$y<L_y$, normal $-\mathbf{e}_y$), replace $\partial/\partial y$ by $-\partial/\partial y$.

\paragraph{Wall-normal velocity at the interface.} As in Wells \& Worster (2011), we use the
reference frame in which the interface is fixed at $y=0$ (it is quasi-planar and recedes slowly). In
this frame the ice moves towards the interface at speed $m$, and melts there. Mass conservation across
the interface, $\rho_i\, m = \rho_w\, v(0)$, gives the boundary condition
\begin{equation}
  v(0) = \frac{\rho_i}{\rho_w}\, m \;>\; 0 , \qquad u(0) = 0 . \label{eq:v0}
\end{equation}
The velocity at the interface is positive: meltwater flows from the interface into the water. It is
the ablation rate times $\rho_i/\rho_w$ ($\approx 0.89$ for sea water, $\approx 0.92$ for fresh
water), i.e.\ it equals the volume flux of meltwater per unit area. It combines the advection of the
water past the receding interface and the expansion of the ice on melting. This is condition (3.5a)
of Wells \& Worster (2011), $v=-(\rho_s/\rho_l)\,\dot a$ with $\dot a=-m$. (In the frame of the ice
the water velocity at the interface is $v(0)-m = (\rho_i/\rho_w-1)\,m<0$: the volume contracts on
melting.)

The interface is therefore \emph{not} impermeable. The heat and salt balances \eqref{eq:heat} and
\eqref{eq:salt} nevertheless contain only the diffusive fluxes on the water side, in either frame. In
the interface-fixed frame the salt balance reads
$\rho_i\, m\, S_i = \rho_w v(0)\, S_b - \rho_w \kappa_S\, \partial S/\partial y|_b$ (salt carried in by
the ice $=$ salt carried away by advection and diffusion); with \eqref{eq:v0} this is exactly
\eqref{eq:salt}. The same holds for heat.

\section{Special cases}

\paragraph{Stefan condition (fresh water).} Without salt, \eqref{eq:salt} drops out and
\eqref{eq:liquidus} becomes $T_b = T_m$ (0\,$^\circ$C):
\begin{equation}
  T_b = T_m, \qquad
  \rho_i L_f\, m = \rho_w c_w \kappa_T \left.\frac{\partial T}{\partial y}\right|_b
  - \rho_i c_i \kappa_i \left.\frac{\partial T_i}{\partial y}\right|_b .
\end{equation}

\paragraph{Transfer-coefficient form (unresolved boundary layers).} Ocean models do not resolve the
diffusive sublayers and write the fluxes with transfer coefficients $\gamma_T$, $\gamma_S$ and the
friction velocity $u_*$:
\begin{equation}
  \kappa_T \left.\frac{\partial T}{\partial y}\right|_b \to \gamma_T u_* (T_w - T_b), \qquad
  \kappa_S \left.\frac{\partial S}{\partial y}\right|_b \to \gamma_S u_* (S_w - S_b),
\end{equation}
with $T_w$, $S_w$ the values in the water outside the sublayer. Neglecting conduction into the ice
and with $S_i=0$, eliminating $m$ and $T_b$ gives a quadratic for $S_b$:
\begin{equation}
  a\, S_b^2 + b\, S_b + c = 0, \qquad
  \begin{aligned}
    a &= c_w \gamma_T \lambda_1, \\
    b &= -\left[c_w \gamma_T (T_w - \lambda_2) + L_f \gamma_S\right], \\
    c &= L_f \gamma_S\, S_w ,
  \end{aligned}
  \label{eq:quadratic}
\end{equation}
Since $\lambda_1<0$, $a<0$ while $c>0$, so the two roots have opposite signs and the physical root is
the positive one (it satisfies $0 < S_b \le S_w$ for melting). Then $T_b$ follows from
\eqref{eq:liquidus} and $m = (\rho_w/\rho_i)\,\gamma_S u_* (S_w - S_b)/S_b$.

In a resolved simulation the gradients are computed directly at the wall and no transfer coefficients
are needed. If the wall gradient is linearised as
$\partial T/\partial y|_b \approx (T_1 - T_b)/\delta_T$ (and similarly for $S$), the same quadratic
\eqref{eq:quadratic} follows with $\gamma_T u_* \to \kappa_T/\delta_T$,
$\gamma_S u_* \to \kappa_S/\delta_S$, $T_w \to T_1$ and $S_w \to S_1$.

\section{Non-dimensional form in \texttt{solver2d}}

\texttt{solver2d} uses a length unit $\ell$, a velocity unit $U$, and $Re = U\ell/\nu$. The
diffusivities are $\kappa_T = U\ell/(Re\,Pr)$ and $\kappa_S = U\ell/(Re\,Sc)$. Write the scalars as
\begin{equation}
  T = T_{\mathrm{ref}} + \Delta T\, \theta, \qquad S = S_{\mathrm{ref}} + \Delta S\, s ,
\end{equation}
($\theta$ and $s$ are scalars 1 and 2 of the code) and the melt rate as a water velocity in units of
$U$:
\begin{equation}
  m^* = \frac{\rho_i\, m}{\rho_w\, U} .
\end{equation}
With $y$ in units of $\ell$, equations \eqref{eq:liquidus}--\eqref{eq:salt} become
\begin{align}
  \theta_b &= \Lambda\, s_b + \theta_{f}, &
  \Lambda &= \frac{\lambda_1 \Delta S}{\Delta T}, \quad
  \theta_f = \frac{\lambda_1 S_{\mathrm{ref}} + \lambda_2 - T_{\mathrm{ref}}}{\Delta T},
  \label{eq:nd-liquidus}\\[0.4em]
  m^* &= St \left( \frac{1}{Re\,Pr} \left.\frac{\partial \theta}{\partial y}\right|_b - q_i \right), &
  St &= \frac{c_w \Delta T}{L_f},
  \label{eq:nd-heat}\\[0.4em]
  \frac{1}{Re\,Sc} \left.\frac{\partial s}{\partial y}\right|_b &= m^*\, (s_b + \sigma), &
  \sigma &= \frac{S_{\mathrm{ref}} - S_i}{\Delta S},
  \label{eq:nd-salt}
\end{align}
with the Stefan number $St$ and the non-dimensional conduction into the ice
\begin{equation}
  q_i = \frac{\rho_i c_i \kappa_i}{\rho_w c_w U \Delta T} \left.\frac{\partial T_i}{\partial y}\right|_b
  \qquad (q_i = 0 \text{ if conduction into the ice is neglected}).
\end{equation}
Equation \eqref{eq:nd-salt} uses $S_b - S_i = \Delta S\,(s_b + \sigma)$. If $S_{\mathrm{ref}}$ is
chosen as the ice salinity ($S_{\mathrm{ref}} = S_i$, e.g.\ $0$), then $\sigma = 0$.

The non-dimensional melt rate is exactly the non-dimensional wall-normal velocity at the interface,
\eqref{eq:v0}:
\begin{equation}
  \frac{v(0)}{U} = m^* > 0 .
\end{equation}

The new non-dimensional parameters are $St$, $\Lambda$, $\theta_f$, $\sigma$ (and $q_i$). In
\texttt{solver2d} the buoyancy enters through the Richardson numbers \texttt{ri(1)}, \texttt{ri(2)}
of $\theta$ and $s$, and the diffusivities through \texttt{pr(1)} $=Pr$ and \texttt{pr(2)} $=Sc$.

\section{Implementation in \texttt{solver2d}}

The model is switched on with \texttt{melt = .true.} in \texttt{\&melt\_model} (source:
\texttt{src/melt.f90}). The parameters are $St$ (\texttt{melt\_st}), $\Lambda$ (\texttt{melt\_lambda}),
$\theta_f$ (\texttt{melt\_thf}) and $\sigma$ (\texttt{melt\_sigma}). The scalars \texttt{melt\_t} and
\texttt{melt\_s} are $\theta$ and $s$; with \texttt{melt\_s = 0} the Stefan condition is used. Heat
conduction into the ice is neglected ($q_i=0$).

\paragraph{Interface values as wall values.} The conditions hold at every point $x$ of the interface,
and are nonlinear ($m^* s_b$) and coupled ($\theta$ and $s$). $\theta_b(x)$ and $s_b(x)$ are imposed as
Dirichlet values on all Fourier modes of the (batched) scalar solves of each Runge--Kutta substep.

\paragraph{Implicit coupling (\texttt{melt\_xdiff = 'implicit'}).} The conditions must hold with the wall gradients of the \emph{new}
scalars. Lagging the interior values by one substep is unstable when diffusion across the first grid
cell within a time step is fast, $\kappa\,\Delta t/\Delta y_1^2 \gtrsim 1$. The theta-method solve is
linear in its wall value $w$: for every Fourier mode $k$ the new wall gradient is
\begin{equation}
  g_k = g^0_k + R_k\, w_k ,
\end{equation}
where $R_k<0$ is the response of the wall gradient to a unit wall value. $R_k$ is computed with the
solver itself (zero right-hand side, unit wall value), for $\theta$ and $s$ separately, and cached per
substep and $\Delta t$. The coupled problem is solved by the iteration
\begin{enumerate}
  \item solve all scalars with the current wall values $w=(\theta_b,s_b)$ and evaluate the wall
        gradients $g$ (first row of the compact $D_1$);
  \item form $G = g - R_{\mathrm{ref}}\, w$ with a constant $R_{\mathrm{ref}}$ per scalar, and solve at
        every $x$ the local conditions with the wall gradient $G + R_{\mathrm{ref}}\, w^{\mathrm{new}}$;
  \item repeat until $\max|w^{\mathrm{new}}-w| < $ \texttt{melt\_tol}.
\end{enumerate}
At convergence the conditions hold exactly with the new gradients ($R_{\mathrm{ref}}$ cancels). For mode
$k$ the error is multiplied by at most $|1-R_k/R_{\mathrm{ref}}|$, which is below 1 for
$R_{\mathrm{ref}} = (\min_k R_k + \max_k R_k)/2$ whatever the time step.

\paragraph{Exact coupling without iteration (default).} The iteration is needed only because $R_k$
depends on $k$, through the implicit $x$-diffusion term $c\,k^2$ of the $\theta$-method, while the
nonlinear local conditions couple the Fourier modes. With \texttt{melt\_xdiff = 'explicit'} the
$x$-diffusion $-\kappa k^2 s$ of $\theta$ and $s$ is advanced explicitly with the RK3 weights, and only
the $y$-diffusion is implicit. The $y$-operator $1 - c\,D_2$ is then the same for every mode, so the
response $R$ and the unit response profile $h(y)$ ($(1-cD_2)h=0$, $h(0)=1$, homogeneous condition at
$y=L_y$) do not depend on $k$, and in physical space
\begin{equation}
  g(x) = g^0(x) + R\, w(x)
\end{equation}
holds at every $x$. The condition is then imposed exactly in one pass per substep: (i) solve with zero
wall values, which gives $s^0$ and $g^0$; (ii) solve the local problem below with $a_T=R_\theta$,
$a_S=R_s$ and $G=g^0$; (iii) set $s = s^0 + w(x)\,h(y)$, i.e.\ $\hat s_k = \hat s^0_k + \hat w_k\, h$ for
every mode. The cost is one scalar solve per substep instead of one per iteration. The explicit
$x$-diffusion adds the RK3 stability limit $\Delta t \le 2.5/(\kappa\, k_{x,\max}^2)$ (with the safety
factor \texttt{cfl\_visc}); since $\Delta x$ is much larger than the near-wall $\Delta y$, this is
usually much weaker than the limits that apply anyway.

\paragraph{Local problem.} With $a_T$, $a_S$ the values of $R_{\mathrm{ref}}$ for $\theta$ and $s$,
$z = s_b+\sigma$ and $G_z = G_s - a_S\sigma$, equations
\eqref{eq:nd-liquidus}--\eqref{eq:nd-salt} reduce at every $x$ to
\begin{equation}
  B z^2 + (A - \kappa_S a_S)\, z - \kappa_S G_z = 0, \qquad
  A = St\,\kappa_T\left[G_\theta + a_T(\theta_f - \Lambda\sigma)\right], \quad
  B = St\,\kappa_T\, a_T \Lambda \ge 0 ,
\end{equation}
with $\kappa_T = 1/(Re\,Pr)$, $\kappa_S = 1/(Re\,Sc)$. Its roots have opposite signs ($a<0$,
$\Lambda\le 0$, $G_z>0$ for positive salinity in the water): the physical root is the non-negative one.
Then $s_b = z-\sigma$, $\theta_b = \Lambda s_b + \theta_f$ and $m^* = A + Bz$. The local problem is solved
in physical space (dealiased grid) and the results are transformed back to Fourier space.

\paragraph{Velocity.} The scalars are solved before the velocity in each substep, and the velocity
solve uses $v(0,x) = m^*(x)$, $u(0,x)=0$ (\texttt{melt\_blowing}), an inhomogeneous wall value of $v$
handled by the same superposition as the other wall conditions. In 2D the mean ($k_x=0$) part of
continuity gives $\partial \bar v/\partial y=0$, so the mean meltwater velocity
$\bar v = \overline{m^*}$ is the same at every $y$ and leaves the domain at $y=L_y$: the whole water
column moves away from the ice at $\overline{m^*}$, as the far field does in the frame of the interface.

\paragraph{Output.} \texttt{log.dat} gives the $x$-averages of $m^*$, $\theta_b$ and $s_b$ and the
number of coupling iterations; \texttt{melt\_NNNNNNN.dat} gives $m^*(x)$, $\theta_b(x)$, $s_b(x)$ at
every field output.

\paragraph{Validation.} Steady melting has exact solutions in the frame of the interface. With the
meltwater velocity $V=m^*$ the profiles are exponential:
$\theta = \theta_b + (\theta_t-\theta_b)(e^{Vy/\kappa_T}-1)/(e^{VL_y/\kappa_T}-1)$ and, since the
total salt flux vanishes, $z = z_b\, e^{Vy/\kappa_S}$. The heat condition then gives
$1 = St\,(\theta_t - \theta_b)/(e^{VL_y/\kappa_T}-1)$; without salt, $V = (\kappa_T/L_y)\ln(1+St)$.
The code reproduces these to the discretisation error (relative $6\times10^{-9}$ for $m^*$ in the
Stefan case; $1.4\times10^{-8}$ and $1.6\times10^{-9}$ on 65 and 129 points for the three-equation case).

\paragraph{Time stepping.} Crank--Nicolson ($\theta$-method with $\theta=1/2$) is not L-stable: after an
abrupt change at the interface, e.g.\ from an initial state that does not satisfy the melt condition, the
wall gradient and hence $m^*$ oscillate from step to step, and the oscillation decays slowly on fine
grids. The implicit Euler method (\texttt{scal\_theta = 1}) damps it and can be used to start a run.

\section{Comparison with Wells \& Worster (2011)}

Wells \& Worster (2011) write the interface as $y=a(t)$, with the solid at $y<a$ and the liquid at
$y>a$, so that $\dot a=\partial a/\partial t<0$ for ablation; subscripts $s$ (solid) and $l$ (liquid),
solute concentration $C$ and solutal diffusivity $D$. With pressure neglected, their interface
conditions are identical to \eqref{eq:liquidus}--\eqref{eq:v0} under the correspondence
\begin{gather*}
  \dot a = -m,\qquad \rho_s=\rho_i,\qquad \rho_l=\rho_w,\qquad c_l=c_w,\qquad \kappa_l=\kappa_T,\\
  D=\kappa_S,\qquad C=S,\qquad C_s=S_i,\qquad T_i^{\mathrm{WW}}=T_b,\qquad C_i=S_b .
\end{gather*}

\begin{center}
\small
\begin{tabular}{l p{7.6cm} p{5.2cm}}
\toprule
 & Wells \& Worster (2011) & this document \\
\midrule
liquidus & $T_i = T_m - \Gamma\,(C_i - C_s)$ \hfill (2.1) & $T_b = \lambda_1 S_b + \lambda_2$:
  \ $\lambda_1=-\Gamma$, $\lambda_2 = T_m + \Gamma C_s$ \\[0.3em]
heat & $\rho_s L \dot a = \rho_s c_s \kappa_s T_y|_{a^-} - \rho_l c_l \kappa_l T_y|_{a^+}$ \hfill (2.2)
  & \eqref{eq:heat} with $\dot a=-m$ \\[0.3em]
salt & $\rho_s (C_i - C_s)\,\dot a = -\rho_l D\, C_y|_{a^+}$ \hfill (2.3) & \eqref{eq:salt} with $\dot a=-m$ \\[0.3em]
velocity & $v(0) = -(\rho_s/\rho_l)\,\dot a$ \hfill (3.5a) & \eqref{eq:v0}: $v(0)=(\rho_i/\rho_w)\,m$ \\
\bottomrule
\end{tabular}
\end{center}

Remarks:
\begin{itemize}
  \item Liquidus: Wells \& Worster use $T_m = 0$\,$^\circ$C and $\Gamma \approx 0.06$\,$^\circ$C\,psu$^{-1}$
        with $C_s\approx 0$, i.e.\ $\lambda_1\approx -0.06$\,$^\circ$C\,psu$^{-1}$ and $\lambda_2 = 0$.
        Holland \& Jenkins (1999) fit $\lambda_1\approx-0.057$\,$^\circ$C\,psu$^{-1}$ and
        $\lambda_2\approx 0.083$\,$^\circ$C over the oceanic salinity range; the difference in $T_b$ is
        below $0.1$\,$^\circ$C.
  \item Heat: Wells \& Worster show (their appendix) that the conductive heat flux into the solid is
        typically small and neglect it, i.e.\ $q_i=0$ in \eqref{eq:nd-heat}.
  \item Velocity: in the interface-fixed frame the wall-normal velocity at $y=0$ equals the meltwater
        volume flux $(\rho_i/\rho_w)\,m$ and is positive (directed into the water). It equals the
        ablation rate $m$ itself only if $\rho_i=\rho_w$.
  \item Their ablation rate, (3.15), $\dot a = -(\rho_l/\rho_s)\,v(0)$, is \eqref{eq:v0} solved for $m$.
        They also note that the fluid flow, and $v(0)$ itself, do not depend on $\rho_s/\rho_l$; only the
        conversion from $v(0)$ to the ablation rate does. This agrees with \eqref{eq:nd-heat}, which
        determines $m^* = v(0)/U$ without $\rho_i/\rho_w$.
  \item Buoyancy: Wells \& Worster keep only the compositional (salt) buoyancy,
        $\rho=\rho_\infty+\rho_l\beta(C-C_\infty)$; \texttt{solver2d} has both, through \texttt{ri(1)}
        and \texttt{ri(2)} (set \texttt{ri(1)} $=0$ to reproduce their assumption).
\end{itemize}

\section*{References}
\begin{itemize}
  \item H.~H. Hellmer and D.~J. Olbers, A two-dimensional model for the thermohaline circulation
        under an ice shelf, \emph{Antarctic Science} 1 (1989) 325--336.
  \item D.~M. Holland and A. Jenkins, Modeling thermodynamic ice--ocean interactions at the base of an
        ice shelf, \emph{J. Phys. Oceanogr.} 29 (1999) 1787--1800.
  \item A.~J. Wells and M.~G. Worster, Melting and dissolving of a vertical solid surface with laminar
        compositional convection, \emph{J. Fluid Mech.} 687 (2011) 118--140.
\end{itemize}

\end{document}
